\section{NN-S015: parity in an infinitesimal Frobenius extension invisible to zeta}

\paragraph{Statement and inputs.}
For every 3-CNF $\Phi$, let $F_\Phi$ be the degree-$r$ form in
$r$ variables of NN-C023, with odd $r\ge3$, and put
\[
 c_\Phi=[z_1\cdots z_r]F_\Phi
       =\#\operatorname{SAT}(\Phi)\pmod2,\qquad
 H=\sum_{i=1}^r z_i^r.
\]
Over the fixed dual-number base
$B=\mathbb F_2[\epsilon]/(\epsilon^2)$, define
\[
 \pi_\Phi:\mathcal X_\Phi
   =V(H+\epsilon F_\Phi)\subset\mathbb P^{r-1}_B
                  \longrightarrow\operatorname{Spec}B.
\]
This is a smooth projective relative hypersurface, constructed in
polynomial time by a compact equation. Its reduction is the fixed
Fermat hypersurface $X_0=V(H)$ for every input with the same $r$.
The based coherent line
\[
 M_\Phi=H^{r-2}(\mathcal X_\Phi,\mathcal O)=B e
\]
has absolute Frobenius given on its underlying two-dimensional
$\mathbb F_2$-space, with ordered basis $(e,\epsilon e)$, by
\[
       N_\Phi=
       \begin{pmatrix}0&0\\c_\Phi&0\end{pmatrix}.       \tag{S015.1}
\]
Thus computing its rank, testing whether it is zero, or deciding whether
the resulting extension of two zero-Frobenius lines splits is
parity-counting complete on this class.

All positive-power traces of $N_\Phi$, its characteristic polynomial,
and its class in the Grothendieck group of finite-dimensional
$\mathbb F_2$-spaces with an endomorphism are independent of $c_\Phi$.
Moreover every finite-field point count and the entire arithmetic zeta
function of $\mathcal X_\Phi$ are identical to those of $X_0$.
For each odd prime $\ell$, their geometric $\ell$-adic cohomology groups
are canonically isomorphic with Frobenius. These are pure because $X_0$
is smooth and projective. Here relative smoothness is over the
nonreduced base $B$; the total $\mathbb F_2$-scheme is not claimed smooth.

The inputs are the coefficient encoding and complete coherent calculation
of NN-C023, its smooth-pencil proof NN-C024, SGA~1, Expos\'e~II,
Theorem~4.10(iii) and Corollary~4.12, and the topological invariance
theorem SGA~4, Expos\'e~VIII, Theorem~1.1, Corollary~1.2 and
Example~1.3(a). The latter have exact locators in the primary re-edition
\url{https://www.normalesup.org/~forgogozo/SGA4/08/08.pdf},
pp.~1--2, retaining original page markers 367--369.
Purity is Deligne, \emph{La conjecture de Weil I}, Theorem~1.6,
printed p.~276. The exact-category Grothendieck-group convention is
SGA~6, Expos\'e~IV, equation~(1.5.2) and Definition~2.2.
No use of A or M is made.

\paragraph{Proof of smoothness and the based coherent line.}
The matrix in NN-C024 that expresses selected multiples of the partial
derivatives in a common monomial degree is $I+uA$. At $u=\epsilon$,
it has the explicit inverse $I-\epsilon A$. Thus all monomials of
that degree lie in the partial-derivative ideal over $B$ itself.
At every projective point some homogeneous coordinate is invertible,
so these partial derivatives generate the unit ideal locally.
Euler's identity, with odd degree $r$, converts this to the
nonvanishing of a relative chart derivative on the hypersurface.
The Jacobian criterion in SGA~1 II~4.10(iii)/4.12 gives smoothness
over $B$. Projectivity follows from the displayed closed immersion.

Multiplication by $H+\epsilon F_\Phi$ is injective on the polynomial
ring over $B$. Indeed writing a putative annihilator as
$A_0+\epsilon A_1$ gives $HA_0=0$ and then $HA_1=0$ in the
polynomial ring over $\mathbb F_2$, so both vanish. This gives
the same exact hypersurface sequence as NN-C023, now over $B$:
\[
 0\longrightarrow\mathcal O_{\mathbb P^{r-1}_B}(-r)
   \xrightarrow{\ H+\epsilon F_\Phi\ }
     \mathcal O_{\mathbb P^{r-1}_B}
   \longrightarrow\mathcal O_{\mathcal X_\Phi}\longrightarrow0.
\]
The coordinate-cover monomial computation works over any coefficient
ring: the sole top Laurent monomial of degree $-r$ is
$1/(z_1\cdots z_r)$, and all intermediate cohomology vanishes.
Hence the boundary map identifies $M_\Phi$ with the free line
on this monomial. It also gives $H^0(\mathcal O_{\mathcal X_\Phi})=B$
and no other nonzero coherent groups. This is an actual based module,
not merely a rank assertion inferred from a special fibre.

The natural Frobenius diagram from NN-C023 applies over $B$.
Absolute Frobenius is semilinear with respect to the square map on $B$;
in the displayed basis it satisfies
\[
       \varphi(b e)=b^2\epsilon c_\Phi e.
\]
The coefficient of $z_1\cdots z_r$ in $H$ is zero, which gives
this formula. For $b=b_0+\epsilon b_1$ one has $b^2=b_0$, with
$b_i\in\mathbb F_2$. Thus $\varphi(e)=\epsilon c_\Phi e$ and
$\varphi(\epsilon e)=0$, proving (S015.1). The linearized relative
Frobenius between the base-pulled coherent line and $M_\Phi$ is
multiplication by $\epsilon c_\Phi$. In this particular example the
underlying absolute map happens also to be $B$-linear, because
$(b^2-b)\epsilon=0$; no such assertion is made for general bases.

\paragraph{Proof of the extension and information-loss statements.}
Let $T$ act as $N_\Phi$, and set $k=\mathbb F_2$.
The invariant subspace $k\epsilon e$ gives an exact sequence of
$k[T]$-modules
\[
 0\longrightarrow(k,T=0)\longrightarrow(M_\Phi,N_\Phi)
     \longrightarrow(k,T=0)\longrightarrow0.          \tag{S015.2}
\]
When $c_\Phi=0$, it splits. When $c_\Phi=1$, any lift of the
quotient basis is $e+a\epsilon e$ and is sent by $T$ to
$\epsilon e\ne0$, so no invariant complementary zero-action
line exists. Thus it does not split.
More explicitly, applying $\operatorname{Hom}_{k[T]}(-,k)$ to
the free resolution
\[
 0\longrightarrow k[T]\xrightarrow{T}k[T]
       \longrightarrow k\longrightarrow0
\]
gives $\operatorname{Ext}^1_{k[T]}(k,k)=k$, since the induced
differential is zero. In these fixed bases, the class of
(S015.2) is exactly $c_\Phi$.

Both matrices in (S015.1) have characteristic polynomial $u^2$,
all positive-power traces zero, and the same associated graded
zero-action object. Relation (S015.2) gives the same class
$2[(k,0)]$ in the exact-category Grothendieck group. Their ranks,
however, are zero and one respectively. Hence trace data and the
additive Grothendieck class cannot reconstruct this extension class.
This is an explicit failure of determination, independent of any
unproved computational lower bound.

The two alternatives already occur with the same $r=3$ in the
encoding: one unconstrained Boolean variable has two assignments
and gives $c_\Phi=0$, while its positive unit clause has one assignment
and gives $c_\Phi=1$. Both have one factor--variable incidence after
the tautological-factor convention of NN-C023. Thus the common
special fibre can be the very same Fermat cubic; changing the
input size is not responsible for the lost information.

There is also an exact orbit-closure interpretation. Over any extension
of $k$, write $N=N_1$. For $t\ne0$,
\[
 \operatorname{diag}(1,t)N\operatorname{diag}(1,t)^{-1}=tN.
\]
The polynomial family $tN$ specializes to zero at $t=0$. Thus zero
lies in the closure of the algebraic $\operatorname{GL}_2$-orbit of
$N$, and every regular polynomial function invariant under this
algebraic group has the same value on the two matrices. Indeed its
value on $tN$ is constant for invertible $t$ and hence, as a polynomial
identity, is the same at zero. This concerns algebraic-group invariants,
not just invariants under the finite group $\operatorname{GL}_2(k)$.
The rank distinction is lost by this invariant quotient as well as by
traces and $K_0$. Computing the rank of an explicitly supplied
two-by-two matrix remains trivial; the parity-complete operation is
recovering its Frobenius action from the compact hypersurface equation.

For a finite extension $K/\mathbb F_2$, every map
$\operatorname{Spec}K\to\mathcal X_\Phi$ sends $\epsilon$ to zero.
It therefore factors uniquely through $X_0$. This proves the
equality of point counts over every finite extension, and hence
of the zeta functions. Equivalently the reduction immersion has
the same closed points and residue fields. This equality is for
the full arithmetic zeta function, not just its reduction modulo two.

The closed immersion $X_0\hookrightarrow\mathcal X_\Phi$ is defined
by a nilpotent ideal. SGA~4 VIII~1.1,1.2,1.3(a) gives the
equivalence of their small \'etale topoi and the corresponding
cohomology isomorphisms. Apply this to each finite
$\mathbb Z/\ell^n$ coefficient sheaf and then pass to the inverse
limit and tensor with $\mathbb Q_\ell$. Naturality makes the
isomorphism compatible with arithmetic Frobenius. The special
fibre $X_0$ is smooth projective because its homogeneous derivatives
are $z_i^{r-1}$; Deligne's theorem gives the asserted purity.
Thus even the entire finite Weil-cohomology system does not
determine the coherent extension in (S015.2).

The reduction $\Phi\mapsto\mathcal X_\Phi$ has polynomial size,
and rank one of $N_\Phi$ is equivalent to odd parity-SAT. Membership
in parity-P on this encoded class is the monomial-choice computation
in NN-C023; complementing the parity convention covers a zero or
splitting test. This proves the completeness assertion.
A deterministic polynomial algorithm for these tasks would give
$\mathrm P=\oplus\mathrm P$ and then $\mathrm{NP}=\mathrm{RP}$
by the same isolation argument as NN-C006. E instead supplies the
subexponential evaluation through the original formula, with the
same input-parameter bounds as NN-C023. No unconditional circuit
lower bound follows from completeness alone.
\hfill$\square$

\paragraph{Grothendieck formulation and the remaining gate.}
The finding is a theorem about a smooth projective first-order
deformation over a length-two local base, its relatively free
coherent cohomology line, and the extension class of its absolute
Frobenius after forgetting to the residue field. The exact
dictionary is the based hypersurface sequence and (S015.2).
The \'etale invariance theorem explains precisely which realization
forgets this infinitesimal information. The Grothendieck-group
relation forgets the nonsplit extension for a different reason:
it imposes additivity along exact sequences.

All fibres in a fixed dimension have the same pure special-fibre
Weil data, yet computing this coherent extension class is
parity-counting complete from the succinct equation. This refutes
the unrestricted inference that a known pure finite Weil object,
its zeta function, or its $K_0$ class supplies an efficient
computation of all related deformation data. It does not identify
the extension with a GCT canonical-basis multiplicity, nor decide
P versus NP or VP versus VNP. Such an identification and an efficient
coefficient-evaluation construction would be additional hypotheses.
